%% kapa: exam=1 round=24 year=2013 period=2 form=A subject=appraisal_law subject_ko=부동산관계법규 seq=7 number=47
%% source: Q-Net 감정평가사 기출문제 (cst00309) artlSeq=5073802 file=extracted/1008/5073802/2013년 제24회 감정평가사 1차시험 기출문제/2교시(A형).hwp
%% conversion: kordoc 4.24.51 -> tools/gen1.py (auto); review=proofread-sonnet; flags=none
\begin{question}{47}
국토의 계획 및 이용에 관한 법령상 바다인 공유수면의 매립목적이 그 매립구역과 이웃하고 있는 농림지역의 내용과 같은 경우, 그 매립준공구역의 용도지역에 관한 설명으로 옳은 것은?
\begin{choices}
  \item 도시ㆍ군관리계획의 입안 및 결정 절차를 거쳐 농림지역으로 지정된다.
  \item 도시ㆍ군관리계획의 입안 및 결정 절차를 거쳐 시가화조정구역으로 지정된다.
  \item 도시ㆍ군관리계획의 입안 및 결정 절차를 거쳐 별도로 지정된다.
  \item 도시ㆍ군관리계획의 입안 및 결정 절차를 거치지 않고 그 매립의 준공인가일부터 농림지역으로 지정된 것으로 본다.
  \item 도시ㆍ군관리계획의 입안 및 결정 절차를 거치지 않고 그 매립의 준공인가일부터 수산자원보호구역으로 지정된 것으로 본다.
\end{choices}
\end{question}
\answer{4}
